\chapter{Modular Arithmetic의 자원 공유와 Pipeline}

\section{설계 목표를 연산식과 자원으로 동시에 적는다}

최종 \rtlfile{unifiedMUL}의 목표는 두 종류의 modular multiplication을 같은
두 DSP48E1과 arithmetic/register path에서 실행하는 것이다.

\begin{itemize}
  \item ML-DSA: $q=8{,}380{,}417$, 23-bit operand path
  \item Narrow primes: R16 Montgomery multiplication, $R=2^{16}$
  \item Narrow modulus: ML-KEM 3,329, HAETAE 64,513, NTRU+ 3,457,
        Falcon 12,289
  \item Latency: six pipeline stages
  \item Throughput: mode가 유지되는 동안 매 cycle 새 입력 수락, \ii=1
\end{itemize}

하나의 \rtlfile{unifiedBUT}은 \rtlfile{unifiedMUL} 두 개를 사용하므로 전체
butterfly arithmetic은 DSP48E1 네 개를 사용한다.

\section{DSP0: 공통 low product}

두 path 모두 먼저 $A$와 $B$의 low 16-bit product를 필요로 한다.

\begin{lstlisting}
(* use_dsp = "yes" *) wire [38:0] c1_mul_lo_w;
assign c1_mul_lo_w = iA[22:0] * iB[15:0];

always @(posedge iSYS_CLK)
  if (iFSM_START) c1_mul_lo_r <= c1_mul_lo_w;
\end{lstlisting}

ML-DSA에서는 23-bit $B$를 low 16-bit와 high 7-bit로 분할한다.
\[
A B=A B_{lo}+2^{16}A B_{hi}.
\]
DSP0의 low product와 DSP1의 high partial product를 정렬해 더한다. Narrow-prime
mode에서는 DSP0 output의 low 32 bit가 $T=A B$다.

\section{Narrow-prime path: R16 Montgomery multiplication}

Narrow mode는 다음 연산을 수행한다.
\[
T=A B,\qquad
u=Tq^{-1}\bmod 2^{16},\qquad
t=\frac{T-uq}{2^{16}}.
\]
$t<0$이면 $q$를 한 번 더해 canonical result를 만든다. 네 modulus는 모두
$q\equiv1\pmod{2^7}$을 만족하므로
\[
q^{-1}=1+2^7c_q\pmod{2^{16}}
\]
로 쓸 수 있다. 따라서 $u$의 low 7 bit는 $T$의 low 7 bit와 같고, high 9 bit는
선택된 shift와 9-bit addition으로 계산할 수 있다.

실제 RTL은 mode별로 필요한 shifted value를 고른 뒤 하나의 9-bit 합에 넣는다.

\begin{lstlisting}
case (mode)
  CTL_KYBER: begin op0 = -m3; op1 = m5; op2 = m6; end
  CTL_HAETAE:begin op0 =  m4;                    end
  CTL_NTRU:  begin op0 =  m2; op1 = m3;
                    op2 =  m5; op3 = m7;          end
  CTL_FALCON:begin op0 = -m2; op1 = -m3;          end
endcase
corr9 = op0 + op1 + op2 + op3;
\end{lstlisting}

여기서 $m_2,\ldots,m_7$은 $T$의 일부 bit를 고정 shift한 9-bit 값이다. 9-bit
wraparound가 modulo 512 연산을 제공한다. 네 modulus마다 넓은 constant multiplier를
병렬로 만든 뒤 고르는 구조를 피하고, mode select와 작은 addition을 공유한다.

\section{DSP1: mode별 곱과 정렬 산술 공유}

두 번째 DSP는 mode에 따라 다른 역할을 한다.

\begin{longtable}{L{29mm}L{57mm}L{75mm}}
\toprule
Mode & DSP1 multiplication & DSP ALU/PREG 동작 \\
\midrule
ML-DSA & $A\times B_{hi}$ & DSP0 low product의 정렬된 high 부분과 덧셈 \\
Narrow prime & $u\times q$ & 정렬된 $T$에서 $u q$를 뺌 \\
\bottomrule
\end{longtable}

최종 synthesis branch는 DSP48E1의 MREG, ALU, PREG를 사용한다. ML-DSA와
Montgomery path의 fabric mux를 multiplier 바로 앞에 깊게 두지 않고 DSP의
A/D/C 입력과 ALUMODE/INMODE로 역할을 바꾼다. simulation branch는 같은 latency의
behavioral multiplication과 register를 사용하므로, 두 branch가 bit-accurate하고
cycle-accurate한지 regression에서 함께 확인해야 한다.

\section{공유 register는 live bit를 기준으로 설계한다}

각 stage의 payload 폭은 두 mode의 필요한 정보를 담을 만큼만 잡는다.

\begin{longtable}{C{15mm}L{52mm}L{94mm}}
\toprule
Stage & ML-DSA에서 보존할 값 & Narrow mode에서 보존할 값 \\
\midrule
C1 & 23$\times$16 low product & $T=A B$ \\
C2 & low product 일부 & $T$와 $c_qT\bmod512$ \\
C3 & quotient 준비용 high/low 조각 & DSP1 subtraction용 정렬된 $T$ \\
C4 & residual과 borrow 정보 & $T-uq$는 DSP1 PREG에 유지 \\
C5 & base-minus-borrow 결과 & negative result에 조건부 $+q$ \\
C6 & ML-DSA canonical correction & C5 narrow result 출력 \\
\bottomrule
\end{longtable}

두 mode가 같은 cycle에 동시에 실행되지 않으므로 payload의 live range가 겹치지
않는 bit는 공유할 수 있다. 다만 mode mux가 wide payload 전체에 생기지 않도록
양쪽에서 실제로 소비되는 low part만 선택하고, 한 mode에서 의미 없는 upper
part는 다른 mode 값을 그대로 연결한다.

\section{최종 correction addition 공유}

C5에서는 ML-DSA의 base-minus-borrow와 narrow Montgomery result의 조건부 $+q$를
하나의 signed 18-bit adder로 표현한다.

\begin{lstlisting}
shared_a = is_dili ? sign_extend(dili_base)
                   : sign_extend(mont_quotient);
shared_b = is_dili ? -borrow
                   : (mont_negative ? mont_q : 0);
shared_sum = shared_a + shared_b;
\end{lstlisting}

공유 가능성은 ``둘 다 덧셈처럼 보인다''에서 나오지 않는다. 두 연산이 같은
pipeline stage에 있고, mode가 transaction 동안 고정되고, 필요한 signed range가
18 bit 안에 들어간다는 세 조건에서 나온다.

\section{ML-DSA path의 폭을 줄이는 방법}

ML-DSA modulus는
\[
2^{23}\equiv 2^{13}-1\pmod q
\]
관계를 갖는다. 최종 RTL은 23-bit operand split으로 얻은 product 조각을 13-bit
radix 위치에 맞춰 재배열하고, 필요한 low subtraction과 high signed addition을
나눠 계산한다. 이는 알려진 modulus relation을 shared datapath에 맞게 배치한
것이며 새로운 reduction algorithm을 주장하는 근거가 아니다.

중요한 RTL 습관은 중간값의 증명된 범위를 사용해 폭을 정하고 assertion으로
남기는 것이다. 예를 들어 signed high base와 final canonical range가 예상 범위를
벗어나면 simulation에서 즉시 error를 발생시킨다. 이런 assertion은 합성에서
제거되므로 production area에는 포함되지 않는다.

\section{산술 timing의 현재 위치}

최종 standalone 150~MHz 최악 경로는 mode register에서 butterfly modular-add
register까지이며, 7 logic levels와 6.522~ns data delay를 보였다. complete BRAM
system 최악 경로도 butterfly input에서 constant-multiply register까지 10 logic
levels, 6.550~ns다. memory-address path가 정리된 뒤 arithmetic이 다시 timing
closure의 중심이 된 것이다.

다음 개선을 시도한다면 mode fanout, constant-add partition, stage boundary를
matched OOC run으로 비교해야 한다. 단순히 DSP를 더 쓰거나 pipeline latency를
늘리기 전에 \ii=1, 네-DSP limit, layer cycle 계약을 함께 보존해야 한다.

\begin{checklistbox}
\begin{itemize}
  \item mode별 수식을 먼저 쓰고 같은 cycle의 공통 연산을 찾는다.
  \item operand와 중간값의 signed range를 증명해 register 폭을 정한다.
  \item DSP 내부 MREG/PREG와 fabric register의 경계를 따로 센다.
  \item mode mux가 DSP/carry 앞 critical path에 붙는지 report로 확인한다.
  \item latency와 initiation interval을 별도 지표로 검증한다.
\end{itemize}
\end{checklistbox}
